\newpage
\phantomsection
\label{prf:self-negation}

\begin{remark*}[Return]
\hyperref[lem:self-negation]{Return to Lemma}
\end{remark*}

\begin{lemma*}[Self Negation]
In any Boolean ring $R$,
\[
-p = p \qquad \text{for all } p \in R.
\]
\end{lemma*}

\begin{proof}
\textbf{Professional Standard Proof.}~\\
By \hyperref[lem:self-additive-inverse]{Self Additive Inverse}, $p + p = 0$.
By the additive inverse axiom~(7), $p + (-p) = 0$.
Since both $p$ and $-p$ satisfy the equation $p + x = 0$, and additive
inverses are unique in any group (a standard group theory fact applied to
the abelian group $(R,+)$), we conclude $-p = p$.
\end{proof}

\begin{proof}
\textbf{Detailed Learning Proof.}~\\

\textbf{Step 1.} Recall the two equations that both hold in $R$.

From Self Additive Inverse (Lemma~\ref{lem:self-additive-inverse}):
\[
p + p = 0. \tag{i}
\]
From the additive inverse axiom~(7) of the ring definition:
\[
p + (-p) = 0. \tag{ii}
\]

\textbf{Step 2.} Apply uniqueness of additive inverses.

The additive group $(R,+,0)$ is a group (abelian, by axiom~(3)). In any
group, if $a + b = 0$ and $a + c = 0$, then:
\[
b = 0 + b = (a + c) + b \overset{\text{assoc}}{=} a + (c + b)
\overset{\text{comm}}{=} a + (b + c) \overset{\text{assoc}}{=} (a+b)+c = 0+c = c.
\]
Apply with $a = p$, $b = p$ (from (i)) and $c = -p$ (from (ii)):
both $p$ and $-p$ are additive inverses of $p$, so $p = -p$.
\end{proof}

\begin{remark*}[Proof structure]
A one-step uniqueness argument. Self Additive Inverse shows $p$ is its
own additive inverse; the ring axiom asserts $-p$ is also an additive
inverse of $p$; uniqueness of inverses in the underlying group forces
$p = -p$. The proof has no calculation content --- it is purely structural.
\end{remark*}

\begin{remark*}[Dependencies]~\\
\begin{itemize}
  \item \hyperref[lem:self-additive-inverse]{Lemma: Self Additive Inverse}
  \item \hyperref[def:ring]{Definition: Ring} (axiom (7), additive inverse law)
\end{itemize}
\end{remark*}
